%% 自動生成（gen.py）
\bktocchap{1-0}{かずを つかうと、なにが いいの？}{「たくさん」より「12こ」の ほうが、よく つたわるね}{7}{1}{1}{10}{J}
\bktocpart{\R{第}{だい}1\R{部}{ぶ}}{かずの はじまり}{16}
\bktocchap{1-1}{いくつ あるかな？}{おさらと リンゴ、どっちが おおい？}{2}{1}{1}{18}{C}
\bktocchap{1-2}{10の まとまりを つくろう}{10ずつ たばねると、\R{大}{おお}きな かずも すぐ わかる}{3}{1}{1}{34}{C}
\bktocchap{1-3}{なにも ないときの「0」}{ものさしの はじまりは、どうして 0なの？}{6}{1}{1}{52}{C}
\bktocchap{1-4}{たし\R{算}{ざん}は「あわせる」「ふえる」}{リンゴ 3こと 2こ、あわせて いくつ？}{2}{1}{1}{66}{C}
\bktocchap{1-5}{ひき\R{算}{ざん}は「のこり」「ちがい」「かわった \R{量}{りょう}」}{きのうより、どれだけ のびた？}{1}{1}{1}{80}{C}
\bktocchap{1-6}{かずの \R{線}{せん}：\R{数直線}{すう|ちょく|せん}}{ものさしの \R{上}{うえ}を、すすんだり もどったり}{6}{1}{1}{94}{C}
\bktocchap{1-7}{\R{大}{おお}きな かずの たし\R{算}{ざん}・ひき\R{算}{ざん}}{36 ＋ 27 は？ 10の まとまりで かんがえよう}{2}{1}{1}{112}{C}
\clearpage
\bktocpart{\R{第}{だい}2\R{部}{ぶ}}{はかる}{130}
\bktocchap{1-8}{ながさを はかろう}{えんぴつ \R{何本}{なん|ぼん}\R{分}{ぶん}？ \R{人}{ひと}に よって ちがったら こまるね}{7}{1}{1}{132}{C}
\bktocchap{1-9}{かさと おもさ}{\R{大}{おお}きい ほうが、おもいとは かぎらない？}{7}{1}{1}{150}{C}
\bktocchap{1-10}{\R{時}{じ}こくと \R{時間}{じ|かん}}{「3\R{時}{じ}」と「3\R{時間}{じ|かん}」は、どこが ちがう？}{1}{1}{1}{166}{C}
\bktocchap{1-11}{\R{目}{め}もりを よもう}{\R{目}{め}もりと \R{目}{め}もりの あいだは、どう よむ？}{7}{1}{1}{184}{C}
\vfill\bkfukidashi{sora}{ながさ、 かさ、 おもさ、 \R{時間}{じ|かん}。 はかる ものは ちがっても、 「1つ\R{分}{ぶん}の いくつ\R{分}{ぶん}」で はかるのは \R{同}{おな}じだよ。}\vspace{10mm}
\clearpage
\bktocpart{\R{第}{だい}3\R{部}{ぶ}}{かけ\R{算}{ざん}と わり\R{算}{ざん}}{200}
\bktocchap{1-12}{かけ\R{算}{ざん}は「\R{同}{おな}じ かずの まとまり」}{おはじきを ならべると、\R{四角}{し|かく}が できるよ}{2}{1}{1}{202}{C}
\bktocchap{1-13}{\R{九九}{く|く}の \R{表}{ひょう}を たんけんしよう}{\R{表}{ひょう}を ななめに \R{見}{み}ると、ふしぎな きまりが…}{5}{1}{1}{218}{C}
\bktocchap{1-14}{わり\R{算}{ざん}は「\R{同}{おな}じ かずずつ \R{分}{わ}ける」}{12こを 3\R{人}{にん}で \R{分}{わ}けると、ひとり いくつ？}{1}{1}{1}{234}{C}
\bktocchap{1-15}{あまりの ある わり\R{算}{ざん}}{100\R{日}{にち}\R{後}{ご}は、\R{何曜日}{なん|よう|び}？}{4}{1}{1}{250}{C}
\bktocchap{1-16}{1より \R{小}{ちい}さい かず：\R{小数}{しょう|すう}と \R{分数}{ぶん|すう}}{ピザを \R{半分}{はん|ぶん}の \R{半分}{はん|ぶん}に すると？}{7}{1}{1}{266}{C}
\vfill\bkfukidashi{mina}{かけ\R{算}{ざん}と わり\R{算}{ざん}は、 「\R{同}{おな}じ かずの まとまり」で つながって いるよ。 \R{九九}{く|く}が つかえると、 わり\R{算}{ざん}も できるね。}\vspace{10mm}
\clearpage
\bktocpart{\R{第}{だい}4\R{部}{ぶ}}{\R{表}{ひょう}・グラフ・かたち}{282}
\bktocchap{1-17}{\R{表}{ひょう}と ぼうグラフ}{いちばん \R{人気}{にん|き}の きゅう\R{食}{しょく}は？ \R{一目}{ひと|め}で わかる}{5}{1}{1}{284}{C}
\bktocchap{1-18}{おれ\R{線}{せん}グラフ：かわり\R{方}{かた}を \R{見}{み}る}{\R{気温}{き|おん}が いちばん はやく \R{上}{あ}がったのは \R{何時}{なん|じ}ごろ？}{1}{1}{1}{304}{C}
\bktocchap{1-19}{かたち：\R{三角}{さん|かく}・\R{四角}{し|かく}・\R{円}{えん}}{タイヤが \R{四角}{し|かく}だったら、どうなる？}{4}{1}{1}{322}{C}
\bktocchap{1-20}{\R{角}{かく}の \R{大}{おお}きさ}{1\R{回転}{かい|てん}は、どれだけ まわった ことに なる？}{4}{1}{1}{340}{C}
\bktocchap{1-21}{\R{広}{ひろ}さを くらべよう：\R{面積}{めん|せき}}{どっちの へやが \R{広}{ひろ}い？ ならべて くらべよう}{2}{1}{1}{358}{C}
\par\vspace{2mm}
\bktocbacklist{\bkbi{こたえ}{374}\bkbi{きごうカード}{383}\bkbi{たんいのひょう}{385}\bkbi{チェックひょう}{389}\bkbi{つぎのステージへ}{392}\bkbi{ふろく}{393}\bkbi{さくいん}{395}}
